% 2011 年山东卷文科第 11 题：长和宽分别相等的两个矩形（正（主）视图、俯视图）。
% 来源 PDF：普通高考/2011/2011山东文.pdf，SHA-256 bf61cf7bb39a3ccfc57735cd2a7ab34387cbfa6bc86a5e7f8246e13292b0e023。
% 原页：第 1 页（A3 三栏）第 2 栏中部，书页码 691。
% 重绘依据：照原页——上下两个全等矩形轮廓，原页实测均 56.68pt × 28.34pt，上下间距 21.25pt；
%   上矩形正下方标「正（主）视图」，下矩形正下方标「俯视图」。无阴影、无虚线、无箭头。
% 坐标取原页 pt：X = x_页 − 566.71，Y = 425.05 − y_页。
\begin{tikzpicture}[x=0.05cm, y=0.05cm, semithick, font=\small]
  % 正（主）视图
  \draw (0,49.59) rectangle (56.68,77.93);
  \node[below, inner sep=3pt] at (28.34,49.59) {正（主）视图};
  % 俯视图
  \draw (0,0) rectangle (56.68,28.34);
  \node[below, inner sep=3pt] at (28.34,0) {俯视图};
\end{tikzpicture}
